\section{Problem Formulation}
\label{sec:formulation}

\subsection{Search Tree and Verifier}

Let $\mathcal{T} = (\mathcal{S}, \mathcal{A}, P, r, s_0)$ be a search tree where
$\mathcal{S}$ is the set of states (partial solutions), $\mathcal{A}$ is the action
set (token or step choices), $P: \mathcal{S} \times \mathcal{A} \to \mathcal{S}$
is the deterministic transition, and $s_0$ is the root. Each leaf $\ell \in \mathcal{L}$
is a complete candidate solution. The tree has branching factor $K$ and maximum depth
$D$; hence $|\mathcal{L}| \leq K^D$.

A \emph{verifier} is a function $V: \mathcal{S} \to [0,1]$ that assigns a quality
score to any state. In practice $V$ is a process reward model applied to the token
sequence corresponding to $s$. We denote by $\hat{V}(s)$ the noisy observation returned
by a single verifier call:
\[
  \hat{V}(s) = V(s) + \varepsilon_s,
  \qquad \varepsilon_s \overset{\text{i.i.d.}}{\sim} \operatorname{SubGaussian}(0, \sigma^2).
\]

\subsection{Dynamic Verifier Allocation}

A search procedure generates a sequence of states $s_1, s_2, \ldots, s_T$ visited
during $T$ search steps. At each step $t$, an \emph{allocation policy} $\pi$ decides
$b_t \in \{0, 1\}$: whether to invoke the verifier at $s_t$. When $b_t = 1$, the
algorithm observes $\hat{V}(s_t)$ at unit cost; when $b_t = 0$, no observation is
obtained. The total verifier budget consumed is $\mathcal{B}(\pi) = \sum_{t=1}^T b_t$.

Let $\hat{v}_t$ denote the algorithm's estimate of $V(s_t)$: either $\hat{V}(s_t)$
if $b_t = 1$, or a \emph{proxy} derived from prior observations if $b_t = 0$. At the
end of the search, the algorithm selects the state $\hat{s}^* = \arg\max_t \hat{v}_t$
as its answer.

\subsection{Oracle Benchmark and Regret}

\begin{definition}[Oracle allocation]
\label{def:oracle}
The oracle $\pi^*$ observes all $V(s_1),\ldots,V(s_T)$ without cost and allocates
verifier calls to maximize the quality of the returned solution. Formally, let
$s^* = \arg\max_t V(s_t)$ be the best state visited. The oracle achieves quality
$V(s^*)$.
\end{definition}

\begin{definition}[Cumulative regret]
\label{def:regret}
The cumulative regret of policy $\pi$ over $T$ steps is
\[
  R(T) \;=\; \sum_{t=1}^{T} \bigl[V(s^*_t) - V(\hat{s}^*_t)\bigr],
\]
where $s^*_t = \arg\max_{i \leq t} V(s_i)$ is the best state found by the oracle
through step $t$, and $\hat{s}^*_t = \arg\max_{i \leq t} \hat{v}_i$ is the policy's
best estimate.
\end{definition}

\begin{remark}
This regret measures the cumulative quality gap relative to an oracle that knows all
true scores. It differs from classical bandit regret in two respects: (i) the
``arms'' are nodes of a tree, not independent, and (ii) the benchmark is a
\emph{retrospective} oracle rather than the Bayes-optimal arm.
\end{remark}

\subsection{Assumptions}
\label{sec:assumptions}

\begin{assumption}[Lipschitz score continuity]
\label{ass:lipschitz}
There exists a constant $L > 0$ such that for any two states $s, s'$ on the same
root-to-leaf path with tree depths $d(s)$ and $d(s')$,
\[
  |V(s) - V(s')| \;\leq\; L \cdot |d(s) - d(s')|.
\]
\end{assumption}

\begin{remark}[Motivation for Assumption~\ref{ass:lipschitz}]
This assumption captures the intuition that the quality of a partial solution changes
smoothly as more tokens are generated. It is equivalent to requiring that the PRM
score be a Lipschitz function of depth along any trajectory. Figure~\ref{fig:lipschitz}
provides empirical support: on simulated PRM trajectories (noise $\sigma{=}0.05$,
consistent with reported PRM score variance from \citet{lightman2023lets}), the estimated
Lipschitz constant $\hat{L}$ stabilizes rapidly and is consistent across depth gaps.
Violations could occur at decision boundaries where the solution becomes clearly correct
or incorrect; in practice these are isolated events that do not materially affect our
bound (see Remark~\ref{rem:robustness}). Section~\ref{sec:gsm8k} validates the assumption
on real Math-Shepherd PRM scores~\citep{wang2024math}.
\end{remark}

\begin{assumption}[Bounded verifier noise]
\label{ass:noise}
The noise $\varepsilon_s$ is $\sigma^2$-sub-Gaussian with $\sigma \leq \sigma_{\max}$
for a known constant $\sigma_{\max}$.
\end{assumption}

\begin{assumption}[Bounded cost ratio]
\label{ass:cost}
Each verifier call costs $c \in [c_{\min}, c_{\max}]$ with $c_{\min} > 0$.
The total compute budget satisfies $\mathcal{B} \leq C_{\max}$ for a known constant.
\end{assumption}

\subsection{Problem Statement}

\begin{problem}
\label{prob:main}
Design an allocation policy $\pi$ that simultaneously achieves:
\begin{enumerate}[(i)]
  \item $R(T) = O(\sqrt{T \log K})$, and
  \item $\mathcal{B}(\pi) = O(\sqrt{T} \log T)$.
\end{enumerate}
\end{problem}

Condition (i) guarantees near-oracle solution quality; condition (ii) guarantees that
verifier calls are sub-linear in the search budget. Together they show that the
computational cost of verification need not scale linearly with search effort.
